\textbf{1. Kiến thức cần nhớ}
\begin{itemize}
\item \textbf{Hình bình hành} là tứ giác có các cạnh đối song song ( \parallel CD$ và  \parallel BC$).
\item \textbf{Tính chất của hình bình hành}: Trong hình bình hành:
\begin{itemize}
\item Các cạnh đối bằng nhau: $\boxed{AB = CD,\ AD = BC}$.
\item Các góc đối bằng nhau: $\boxed{\widehat{A} = \widehat{C},\ \widehat{B} = \widehat{D}}$.
\item Hai đường chéo cắt nhau tại trung điểm của mỗi đường: $\boxed{OA = OC,\ OB = OD}$.
\end{itemize}
\item \textbf{Dấu hiệu nhận biết hình bình hành} ($ dấu hiệu):
\begin{enumerate}
\item Tứ giác có các cạnh đối song song là hình bình hành.
\item Tứ giác có các cạnh đối bằng nhau là hình bình hành.
\item Tứ giác có \textbf{một cặp cạnh đối song song và bằng nhau} là hình bình hành.
\item Tứ giác có các góc đối bằng nhau là hình bình hành.
\item Tứ giác có hai đường chéo cắt nhau tại trung điểm của mỗi đường là hình bình hành.
\end{enumerate}
\end{itemize}

\textbf{2. Phương pháp giải các dạng}

%% DANG 8H3?12-1
\textbf{Dạng 1. Tính chất hình bình hành}
\begin{enumerate}
\item Vận dụng các cạnh đối song song và bằng nhau để tính độ dài đoạn thẳng.
\item Hai góc kề một cạnh bù nhau ($\widehat{A} + \widehat{B} = 180^\circ$), hai góc đối bằng nhau ($\widehat{A} = \widehat{C}$) để tính số đo góc.
\item Vận dụng giao điểm hai đường chéo là trung điểm mỗi đường để chứng minh các đoạn thẳng bằng nhau.
\end{enumerate}

%% DANG 8H3?12-2
\textbf{Dạng 2. Dấu hiệu nhận biết hình bình hành}
\begin{enumerate}
\item Chọn dấu hiệu phù hợp nhất với giả thiết bài toán:
\begin{itemize}
\item Nếu đề cho các cặp cạnh song song $\rightarrow$ dùng dấu hiệu 1.
\item Nếu đề cho các trung điểm của đoạn thẳng $\rightarrow$ ưu tiên dùng dấu hiệu 5 (hai đường chéo cắt nhau tại trung điểm mỗi đường).
\item Nếu đề cho một cặp cạnh vừa song song vừa bằng nhau $\rightarrow$ dùng dấu hiệu 3 (dấu hiệu hay dùng nhất).
\end{itemize}
\end{enumerate}

%% DANG 8H3?12-3
\textbf{Dạng 3. Chứng minh đoạn thẳng, góc bằng nhau, ba điểm thẳng hàng}
\begin{enumerate}
\item \textbf{Chứng minh hai đoạn thẳng hoặc hai góc bằng nhau}: Chứng minh chúng là hai cạnh đối hoặc hai góc đối của một hình bình hành.
\item \textbf{Chứng minh ba điểm thẳng hàng}: Chứng minh điểm thứ ba là giao điểm của hai đường chéo trong một hình bình hành (hoặc qua một điểm có hai đường thẳng cùng song song với một đường thẳng thứ ba).
\item \textbf{Chứng minh hai đường thẳng song song}: Chứng minh chúng chứa hai cạnh đối của một hình bình hành.
\end{enumerate}
